% page 24 of the facsimile (image p-023 of the scan)
% block 157.5 x 246.3 mm ; left margin 8.0 mm
\pgc{-12.700mm}{0.000mm}{
\l{}
\l{}
\l{}
\l{}
\l{}
\l{\sou{alaktar}. (\sou{trans,}) (\sou{aludante muliero}). Nutrar (infanteto) per}
\l{\cel{3}sua lakto. - EFIS.}
\l{}
\l{\sou{alambiko}.\cel{2}Aparato por distilar. - DEFIRS.}
\l{}
\l{\sou{alarmar}.\cel{2}(\sou{trans}.) Advokar a lia armi korpo de soldati.}
\l{\cel{3}- DEFIS.}
\l{}
\l{\sou{alaudo}. (\sou{zool.}) Ucelo di la familio \textquotedbl{}konirostri\textquotedbl{}, qua vivas}
\l{\cel{3}en la agri. - L.\cel{1}\sou{alauda}. - FISL.}
\l{}
\l{\sou{albo}. (\sou{liturgio katolika}). Tuniko ek telo blanka, kunligita}
\l{\cel{3}super la reni per zono o kordono, quan la episkopi, la}
\l{\cel{3}sacerdoti, e c. +weras sur sua sutano por celebrar la}
\l{\cel{3}meso. - DEFIS.}
\l{}
\l{\sou{albatroso}.\cel{2}(\sou{zool.}) Mar-ucelo di la mi-sfero australa di Tero,}
\l{\cel{3}ek la familio \textquotedbl{}longa-pluma\textquotedbl{}, di qua la beko esas forta e}
\l{\cel{3}hokatra, e la pedi, sen-polexa. - L.\cel{1}\sou{diomedea}. - DEFIRS.}
\l{}
\l{\sou{albergo}. Hotelo, situita (generale) en la pre-urbi di urbo,}
\l{\cel{3}an ruro-voyo od en vilajo. - DFIS.}
\l{}
\l{\sou{albina}. (\sou{patol}.) Anomaleso korpala quan karakterizas, che la}
\l{\cel{3}homo, la senkoloreso di la pelo, di la hari, di la pili,}
\l{\cel{3}la paleso rozea di la iriso, la redatreso di la pupilo,}
\l{\cel{3}e quan onu vidas anke che ula animali (kunikli furo-blanka)}
\l{\cel{3}od elefanti pelo-blanka. - DEFIRS.}
\l{}
\l{\sou{albumo}. Mikra registro, di qua la folii esas destinita a re-}
\l{\cel{3}cevar noturi, autografi, skisuri, versi, muzikaji, e c.}
\l{\cel{3}- DEFIRS.}
\l{}
\l{\sou{albumeno}. Substanco qua envelopas la embriono en ula grani e}
\l{\cel{3}nutras lu dum la jermifo. (\sou{anke pri la ovi)}.- DEFIRS.}
\l{}
\l{\sou{albumino}. Substanco quan onu renkontras en la sero di sango,}
\l{\cel{3}en la albumeno di la ovo, e qua koagulas efike da kaloro.}
\l{\cel{3}- DEFIRS.}
\l{}
\l{\sou{albuminoido}. To di qua la naturo esas olta di albumino; dic-}
\l{\cel{3}esas pri la substanci azotoza (nitroza) di animali o vej-}
\l{\cel{3}etanti. - DEFIRS.}
\l{}
\l{\sou{albuminurio}. (\sou{patol.)}\cel{1}Formo di maladeso dum qua urino esas}
\l{\cel{3}albuminoza. - DEFIRS.}
\l{}
\l{\sou{alburno}. (\sou{bot.}) En la arbori bikotiledona, la strato unesma}
\l{\cel{3}de ligno, ordinare blanka, qua jacas nemediate sub la}
\l{\cel{3}kortico. - EFISL.}
\l{}
\l{\sou{alcedo}. (\sou{zool.}) Mar-hirundo. - L.\cel{1}\sou{alcedo}.}
}
